\section{Mathematics}

This thesis uses standard mathematical notation for rigid-body dynamics and
control theory. Key symbols and conventions are defined here for reference.

\subsection{Coordinate Frames}

The drone's pose is expressed in two coordinate frames:
\begin{itemize}
    \item \textbf{World frame (ENU):} $X$ east, $Y$ north, $Z$ up.
    \item \textbf{Body frame:} $x_b$ forward, $y_b$ right, $z_b$ down
    (NED convention per PX4 firmware).
\end{itemize}

Position is denoted $\mathbf{p} = [x, y, z]^\mathsf{T}$,
velocity $\mathbf{v} = [v_x, v_y, v_z]^\mathsf{T}$, and
acceleration $\mathbf{a} = [a_x, a_y, a_z]^\mathsf{T}$.

\subsection{Drone Dynamics}

The quadcopter's rotational dynamics are governed by the Euler equations:

\begin{equation}
    \mathbf{I} \dot{\boldsymbol{\omega}} + \boldsymbol{\omega} \times (\mathbf{I} \boldsymbol{\omega}) = \boldsymbol{\tau}
\end{equation}

where $\mathbf{I}$ is the inertia tensor, $\boldsymbol{\omega}$ the angular velocity,
and $\boldsymbol{\tau}$ the applied torque from motor thrust differentials.

\subsection{Collision Kinematics}

Impact deceleration is measured as the peak magnitude of the IMU acceleration vector
during the collision window:

\begin{equation}
    \|\mathbf{a}_{\text{peak}}\| = \max_{t \in T_{\text{collision}}} \sqrt{a_x(t)^2 + a_y(t)^2 + a_z(t)^2}
\end{equation}

This is reported in multiples of standard gravity $g = 9.81\ \text{m}\,\text{s}^{-2}$.
